\section{Affine Cipher}
\label{sec:affine}

Affine cipher adalah perluasan langsung dari Caesar cipher. Bila Caesar cipher
hanya menggeser setiap huruf dengan sejumlah langkah yang tetap, affine cipher
melakukan dua operasi sekaligus: \emph{mengalikan} posisi huruf dengan sebuah
bilangan, lalu \emph{menambahkan} sebuah pergeseran. Namanya berasal dari
istilah \textit{affine transformation} di dalam geometri, yaitu pemetaan yang
mempertahankan garis lurus, karena bentuk fungsinya memang fungsi linier
$f(x) = mx + b$. Penambahan parameter perkalian itulah yang memperbesar ruang
kunci dari 25 kemungkinan pada Caesar cipher menjadi 300 kemungkinan pada
affine cipher, meskipun keduanya tetap merupakan cipher abjad-tunggal
\citep{stinson2018}.

\subsection{Rumusan dan Syarat Kunci}
\label{sec:affine-rumusan}

Secara matematis, enkripsi dan dekripsi affine cipher dinyatakan sebagai
\begin{equation}
    C \equiv mP + b \pmod{n},
    \label{eq:affine-enkripsi}
\end{equation}
\begin{equation}
    P \equiv m^{-1}(C - b) \pmod{n},
    \label{eq:affine-dekripsi}
\end{equation}
dengan keterangan sebagai berikut.
\begin{enumerate}
    \item $n$ adalah ukuran alfabet. Untuk alfabet Latin baku, $n = 26$.
    \item $m$ adalah bilangan bulat yang \emph{relatif prima} dengan $n$,
          yaitu $\gcd(m, n) = 1$. Faktor $m$ inilah yang membedakan affine
          cipher dari Caesar cipher.
    \item $b$ adalah jumlah pergeseran, yang peranannya sama seperti kunci pada
          Caesar cipher.
    \item $m^{-1}$ adalah invers perkalian dari $m$ modulo $n$, yaitu bilangan
          yang memenuhi $m \cdot m^{-1} \equiv 1 \pmod{n}$.
    \item Kuncinya adalah pasangan $(m, b)$. Karena itu penulisan kunci affine
          selalu memuat dua bilangan, bukan satu.
    \item Caesar cipher adalah kasus khusus affine cipher dengan $m = 1$;
          dalam hal ini Persamaan~\ref{eq:affine-enkripsi} menyusut menjadi
          $C \equiv P + b \pmod{26}$.
\end{enumerate}

\textbf{Mengapa $m$ harus relatif prima dengan $n$?} Syarat itu bukan sekadar
keindahan matematis, melainkan keharusan agar cipher dapat didekripsi. Fungsi
enkripsi harus merupakan pemetaan satu-ke-satu dari himpunan huruf ke dirinya
sendiri; bila dua huruf plainteks yang berbeda menghasilkan huruf cipherteks
yang sama, maka cipherteks itu tidak dapat dikembalikan secara unik. Sekarang
perhatikan apa yang terjadi bila $\gcd(m, 26) = d > 1$. Karena $m$ dan 26
mempunyai faktor bersama $d$, nilai $mP \bmod 26$ hanya dapat menempuh $26/d$
nilai yang berbeda, sehingga sebagian huruf cipherteks tidak akan pernah
muncul. Sebagai contoh, untuk $m = 2$ dan $b = 0$ berlaku $C \equiv 2P
\pmod{26}$, sehingga hanya huruf cipherteks bernomor genap yang mungkin
muncul, sedangkan
\[
P = 0 \ (\code{A}) \quad\text{dan}\quad P = 13 \ (\code{N})
\]
sama-sama menghasilkan $C = 0$ (\code{A}). Begitu pula $m = 13$ memberikan
$P = 0$ dan $P = 2$ yang keduanya menghasilkan \code{A}. Karena itu $m$ harus
dipilih dari himpunan bilangan yang relatif prima dengan 26, yaitu
\[
m \in \{1, 3, 5, 7, 9, 11, 15, 17, 19, 21, 23, 25\},
\]
yang berjumlah $\varphi(26) = 12$ buah. Perhatikan bahwa himpunan itu memuat
$m = 1$ (Caesar cipher) dan $m = 25$, yang karena $25 \equiv -1 \pmod{26}$
ekuivalen dengan pembalikan urutan alfabet.

Dengan 12 pilihan untuk $m$ dan 25 pilihan untuk $b$ (yaitu pergeseran 1 sampai
25, sebab $b = 0$ berarti tidak ada pergeseran sama sekali), ruang kunci affine
cipher berisi
\[
12 \times 25 = 300
\]
kemungkinan pasangan kunci. Angka itu masih sangat kecil: sebuah komputer
biasa dapat mencoba seluruh 300 kunci itu dalam waktu jauh kurang dari satu
detik. Perluasan ruang kunci memang terjadi, tetapi tidak cukup besar untuk
memberi keamanan. Tabel~\ref{tab:ruang-kunci-abjad-tunggal} membandingkan
ketiga cipher abjad-tunggal yang telah dibahas.

\begin{table}[htbp]
\centering
\small
\begin{tabularx}{\linewidth}{@{}>{\raggedright\arraybackslash}p{3.4cm} r >{\raggedright\arraybackslash}X@{}}
\hline
\textbf{Cipher} & \textbf{Ukuran ruang kunci} & \textbf{Keterangan} \\
\hline
Caesar cipher & 25 & Hanya pergeseran tetap; $m = 1$ dan $b \in \{1, \dots, 25\}$ \\[4pt]
Affine cipher & 300 & $12$ nilai $m$ yang relatif prima dengan 26 dikali 25 nilai
$b$; masih dapat dicoba seluruhnya dalam sekejap \\[4pt]
Substitusi abjad-tunggal & $26! \approx 4{,}03 \times 10^{26}$ & Ruang kunci
raksasa, tetapi tetap runtuh karena distribusi frekuensi huruf tidak berubah
\\
\hline
\end{tabularx}
\caption{Perbandingan ruang kunci tiga cipher abjad-tunggal. Ukuran ruang kunci
yang besar tidak menjamin keamanan bila struktur statistik plainteks masih
terbaca pada cipherteks}
\label{tab:ruang-kunci-abjad-tunggal}
\sumber{Dihitung dari Munir, \textit{IF4020 Kriptografi}, ITB; tabel disusun dari uraian tersebut.}
\end{table}

\subsection{Contoh Enkripsi dan Dekripsi}
\label{sec:affine-contoh}

Misalkan kuncinya adalah $m = 7$ dan $b = 10$, sedangkan plainteksnya adalah
\code{kripto}. Setiap huruf diubah lebih dahulu menjadi bilangan menurut
pemetaan $\code{A} = 0$ sampai $\code{Z} = 25$, kemudian dikenai
Persamaan~\ref{eq:affine-enkripsi}. Tabel~\ref{tab:affine-enkripsi}
memperlihatkan perhitungannya huruf demi huruf.

\begin{table}[htbp]
\centering
\small
\begin{tabular}{@{}c c c c c@{}}
\hline
\textbf{Plainteks} & \textbf{Nilai $P$} & \textbf{$7P + 10$} &
\textbf{$(7P+10) \bmod 26$} & \textbf{Cipherteks} \\
\hline
\code{k} & 10 & 80  & 2  & \code{C} \\[2pt]
\code{r} & 17 & 129 & 25 & \code{Z} \\[2pt]
\code{i} & 8  & 66  & 14 & \code{O} \\[2pt]
\code{p} & 15 & 115 & 11 & \code{L} \\[2pt]
\code{t} & 19 & 143 & 13 & \code{N} \\[2pt]
\code{o} & 14 & 108 & 4  & \code{E} \\
\hline
\end{tabular}
\caption{Enkripsi \code{kripto} dengan affine cipher berkunci $m = 7$ dan
$b = 10$. Hasil akhirnya adalah \code{CZOLNE}}
\label{tab:affine-enkripsi}
\sumber{Data dari Munir, \textit{IF4020 Kriptografi}, ITB; tabel disusun dari contoh tersebut.}
\end{table}

Perhatikan huruf pertama: $\code{k} = 10$ sehingga $7 \cdot 10 + 10 = 80$,
dan $80 \bmod 26 = 2$ karena $80 = 3 \cdot 26 + 2$. Bilangan 2 bersesuaian
dengan huruf \code{C}. Cara yang sama menghasilkan cipherteks
\begin{center}
\texttt{CZOLNE}.
\end{center}

Untuk mendekripsi, diperlukan invers dari $m = 7$ modulo 26. Karena
$7 \cdot 15 = 105 = 4 \cdot 26 + 1$, maka $7^{-1} \equiv 15 \pmod{26}$.
Dengan Persamaan~\ref{eq:affine-dekripsi}, huruf pertama cipherteks
\code{C} $= 2$ memberi
\[
P \equiv 15\,(2 - 10) = 15 \cdot (-8) = -120 \equiv 10 \pmod{26},
\]
dan $10$ bersesuaian dengan huruf \code{k} --- kembali ke plainteks semula.
Sebagai catatan, perhitungan modulo atas bilangan negatif dilakukan dengan
menambahkan kelipatan 26 secukupnya: $-120 + 5 \cdot 26 = 10$. Pembaca yang
belum terbiasa dengan aritmetika modulo sebaiknya meninjau kembali
Bagian~\ref{sec:kekongruenan} pada Bab~\ref{chap:landasan-matematika}.

\subsection{Kriptanalisis Affine Cipher}
\label{sec:kriptanalisis-affine}

Meskipun ruang kuncinya lebih besar daripada Caesar cipher, affine cipher jauh
lebih rapuh daripada yang tersirat oleh angka 300. Kelemahannya bersumber dari
sifatnya yang \emph{linier}: hubungan antara plainteks dan cipherteks dapat
dirumuskan sebagai sistem persamaan linier, dan sistem persamaan linier dapat
diselesaikan secara langsung tanpa pencarian.

\textbf{Serangan dengan known-plaintext.} Misalkan kriptanalis mengetahui dua
huruf plainteks beserta padanan cipherteksnya, yaitu $P_1 \to C_1$ dan
$P_2 \to C_2$. Dari Persamaan~\ref{eq:affine-enkripsi} diperoleh dua
kekongruenan simultan
\begin{align}
    C_1 &\equiv mP_1 + b \pmod{26}, \label{eq:affine-simultan-1} \\
    C_2 &\equiv mP_2 + b \pmod{26}. \label{eq:affine-simultan-2}
\end{align}
Dengan mengurangkan Persamaan~\ref{eq:affine-simultan-2} dari
Persamaan~\ref{eq:affine-simultan-1}, parameter $b$ tereliminasi dan tersisa
\[
C_1 - C_2 \equiv m\,(P_1 - P_2) \pmod{26}.
\]
Bila selisih $P_1 - P_2$ relatif prima dengan 26, maka
\[
m \equiv (C_1 - C_2)\,(P_1 - P_2)^{-1} \pmod{26},
\]
dan setelah $m$ diketahui, nilai $b$ langsung diperoleh dari salah satu
persamaan di atas. Jadi kunci affine cipher dapat dihitung dari dua huruf saja.

Sebagai contoh, misalkan kriptanalis mengamati bahwa huruf \code{C}
berpadanan dengan \code{K} dan huruf \code{E} berpadanan dengan \code{O}.
Dengan nilai numeriknya, $\code{C} = 2$, $\code{K} = 10$, $\code{E} = 4$, dan
$\code{O} = 14$, diperoleh
\[
\begin{aligned}
2 &\equiv 10m + b \pmod{26},\\
4 &\equiv 14m + b \pmod{26}.
\end{aligned}
\]
Pengurangan kedua persamaan menghasilkan
\[
2 \equiv 4m \pmod{26}.
\]
Kekongruenan ini mempunyai dua penyelesaian di dalam rentang $0$ sampai $25$,
yaitu $m = 7$ dan $m = 20$, karena keduanya memberi $4m \bmod 26 = 2$. Namun
$\gcd(20, 26) = 2 \ne 1$, sehingga $m = 20$ tidak sah sebagai kunci affine.
Satu-satunya nilai yang sah adalah $m = 7$. Substitusi ke persamaan pertama
memberi $b \equiv 2 - 70 = -68 \equiv 10 \pmod{26}$, sehingga kunci yang
diperoleh adalah $(m, b) = (7, 10)$ --- persis kunci yang dipakai pada
Tabel~\ref{tab:affine-enkripsi}. Perhatikan bahwa ketika
$\gcd(P_1 - P_2, 26) \ne 1$, penyelesaian $m$ tidak tunggal; kriptanalis
kemudian menyaring calon-calon itu dengan syarat $\gcd(m, 26) = 1$ dan menguji
sisanya pada huruf ketiga.

\textbf{Serangan tanpa known-plaintext.} Bila kriptanalis tidak mempunyai
pasangan plainteks dan cipherteks sama sekali, affine cipher masih dapat
dipecahkan dengan \emph{exhaustive key search}: cukup mencoba seluruh
$12 \times 25 = 300$ pasangan $(m, b)$ dan memeriksa mana yang menghasilkan
teks bermakna. Jumlah 300 itu terlalu kecil untuk menjadi penghalang, sehingga
affine cipher pada praktiknya hanya bertahan bila penyerang tidak berusaha
sama sekali.

\subsubsection{Varian: Affine Cipher atas Blok Huruf}

Salah satu cara memperbesar faktor kerja pencarian menyeluruh adalah dengan
tidak mengenkripsi huruf secara individual, melainkan mengenkripsi blok yang
terdiri atas beberapa huruf sekaligus. Misalkan pesan \code{kriptografi}
dipecah menjadi kelompok empat huruf,
\begin{center}
\texttt{krip togr afi},
\end{center}
dan setiap kelompok dipandang sebagai bilangan desimal delapan angka dengan
setiap hurufnya disandikan sebagai dua angka (karena $A = 0$ sampai $Z = 25$
selalu dapat dinyatakan dengan dua angka). Kelompok pertama menjadi
\[
\code{krip} \rightarrow 10\,17\,08\,15 \rightarrow 10170815,
\]
kelompok kedua \code{togr} $\rightarrow 19140617$, dan kelompok ketiga
\code{afi} $\rightarrow 00\,05\,08 \rightarrow 000508$, dengan angka nol di
depan sebagai pengisi agar panjangnya tetap delapan angka.

Karena bloknya dinyatakan sebagai bilangan delapan angka, modulus yang dipakai
tidak lagi 26. Nilai terbesar yang mungkin muncul adalah $25252525$ (padanan
dari \code{ZZZZ}, yaitu empat huruf bernomor 25 berturut-turut), sehingga
bilangan itu dapat dipilih sebagai modulus $n$. Dengan memilih $m = 21035433$
yang relatif prima dengan $25252525$ dan $b = 23210025$, fungsi enkripsinya
menjadi
\[
C \equiv 21\,035\,433\,P + 23\,210\,025 \pmod{25\,252\,525}.
\]
Sebagai contoh, blok pertama $P = 10170815$ menghasilkan
\[
C = (21\,035\,433 \cdot 10\,170\,815 + 23\,210\,025) \bmod 25\,252\,525
= 22\,837\,395.
\]
Fungsi dekripsinya adalah
\[
P \equiv 11\,837\,547\,(C - 23\,210\,025) \pmod{25\,252\,525},
\]
karena $21\,035\,433 \cdot 11\,837\,547 \equiv 1 \pmod{25\,252\,525}$.
Sebagai pemeriksaan, dekripsi $C = 22\,837\,395$ mengembalikan
$P = 10\,170\,815$, yaitu \code{krip} semula.

Varian blok ini memperbesar ruang kunci menjadi kelipatan $25\,252\,525$,
jauh lebih besar daripada 300. Namun perlu ditekankan bahwa memperbesar ruang
kunci \emph{tidak} memperbaiki kelemahan struktur: affine cipher atas blok
huruf tetap merupakan fungsi linier, sehingga penyerang yang mengetahui dua
pasang blok plainteks--cipherteks tetap dapat menghitung $m$ dan $b$ secara
langsung dengan cara yang sama seperti pada Persamaan~\ref{eq:affine-simultan-1}
dan~\ref{eq:affine-simultan-2}. Pelajaran inilah yang membawa kita ke Hill
cipher pada Bagian~\ref{sec:hill}: cipher blok yang lebih kuat justru karena
memanfaatkan percampuran antarhuruf di dalam blok, bukan sekadar karena
memperbesar alfabet \citep{stinson2018}.
